\documentclass[11pt]{mpreport}
\usepackage[utf8]{inputenc}
\usepackage[T1]{fontenc}
\usepackage[scaled=0.85]{helvet}
\usepackage[colorlinks=true,linkcolor=black,citecolor=black, urlcolor=black]{hyperref}
\usepackage{microtype}
\usepackage{fancyhdr}
\usepackage{graphicx}
\usepackage{palatino}
\usepackage{listings}
\usepackage{appendix}
\usepackage{remreset}
\usepackage{qtree}
\usepackage{wrapfig}
\usepackage{array}
\usepackage{multirow}
\usepackage{amsmath}
\usepackage{amssymb}
\usepackage{color}
\makeatletter
\@removefromreset{footnote}{chapter}
\makeatother



\setlength{\headheight}{15.2pt}
\pagestyle{fancy}
\fancyhead[RO]{\thepage}
\fancyhead[LO]{}
\fancyfoot{}

\newcommand{\TODO}[1]{\hspace{-1.45cm}\textcolor{red}{ \textbf{TODO: #1}}}
\newcommand{\ITODO}[1]{\textcolor{red}{ \textbf{TODO: #1}}}





\definecolor{orange}{rgb}{1,.49,.18}
\definecolor{darkgreen}{rgb}{0,.66,.0}
\definecolor{mypurple}{rgb}{.76,0,.83}
\newcommand{\var}[1]{{\color{orange} \textbf{#1}}}
\newcommand{\autovar}[1]{{\color{darkgreen} \textbf{#1}}}
\newcommand{\conj}[1]{{\color{mypurple} \textbf{#1}}}



\fancypagestyle{plain}{%
   \fancyhf{}%
   \fancyhead[C]{} %Kapitelname ausblenden

   \renewcommand{\headrulewidth}{0.0pt} %obere Linie ausblenden
   \fancyfoot[C]{\thepage}
}

\setcounter{secnumdepth}{2} 


\begin{document}
\lstset{language=Java, basicstyle=\footnotesize}

\renewcommand{\thepage}{\roman{page}}
\newcommand{\monthword}[1]{\ifcase#1\or January\or February\or March\or April\or
                                        May\or June\or July\or August\or
                                        September\or October\or November\or December\fi} 

  \begin{titlepage}
    \def\fakultaet{Faculty of Mathematics and Computer Science}
    \par
    \unitlength1mm
    \begin{picture}(-2,12)
      \put(-1,11){\rule{165mm}{1mm}}
      \put(-4,-17){
        \parbox[b]{130mm}{
          \raggedleft{
            \fontsize{14}{18}\selectfont
            SAARLAND UNIVERSITY\\
            ~\\
            \fontfamily{ptm}\fontsize{12}{14}\selectfont \fakultaet\\
           Department of Computer Science\\
            Related Work Summary\\
          }
        }
      }
      \put(130,-20){\includegraphics[height=30mm]{uds}}
      \put(-1,-20){\rule{130mm}{.25mm}}
    \end{picture}
    
    
    \vspace{6cm}
	\hspace{2cm}\parbox[b]{150mm}{
    \begin{flushleft}
      \sffamily\Huge
	Understanding Autonomous Device Behavior in a Social Internet of Things environment though Personality Traits\\
~\\
     
      
    \end{flushleft}

}
      \vfill
\hspace{3cm}\parbox[b]{120mm}{
      \sffamily \Huge
      Latifa Abdullayeva\\
      \Large
     \monthword{\month} \the\year
}

    \hspace{-5cm}
    \rule{8cm}{2.25mm}
    \hspace{-5cm}
    \rule{18cm}{.25mm}
  \end{titlepage}



\pagestyle{empty}

\vspace*{0.5cm}
\textbf{Advisor:}\\
Donald Degraen\\
German Research Center for Artificial Intelligence\\
Saarland Informatics Campus\\
Saarbrücken, Germany

\vspace{4.5cm}
Started on 01/01/1970\\
Handed in on 08/07/2019


\vspace{4.5cm}


\vspace{3cm}
Saarland University\\
Faculty of Natural Sciences and Technology I\\
Department of Computer Science\\
Campus - Building E1.1\\
66123 Saarbrücken\\
Germany\\




\clearpage

\begingroup
  \pagestyle{plain}

\setcounter{tocdepth}{2}
  \addtocontents{toc}{\protect\thispagestyle{plain}}
  \addtocontents{toc}{\protect\vspace{-.45cm}}
  \tableofcontents
  \clearpage
\endgroup 

\pagestyle{fancy}
\setcounter{page}{1}
\renewcommand{\thepage}{\arabic{page}}

\section{Related work}
This chapter presents the background and current works related to our research and is organised as follows. Sections ~\ref{sec:Ubiquitous Computing}, ~\ref{sec:Internet of Things}, and ~\ref{sec:Social Internet of Things} provide an introduction of basic concepts related to our prototype. Then, in section ~\ref{sec:Autonomous interaction of things}, we present a system that inspired us to expand it. Starting from section ~\ref{sec:The Theory of Proxemics}, we discuss some methodologies that help us to expand the existing system. 

\subsection{Ubiquitous Computing}
\label{sec:Ubiquitous Computing}
The attempts to make technologies invisible in the background of people’s life led to the emergence of a new approach in the area of Information Technology whereby making the term Ubiquitous Computing prominent in recent years. The “Ubiquitous Computing” was initially put forward by Mark Weiser in “The computer for the 21st Century” paper ~\cite{weiser2002computer}. In this paper, the author touched two issues related to the concept of Ubiquitous Computing such as the location and scale.
\par The traditional computers which existed before the introduction of this paradigm had no idea about their location. The location-aware system may have information about how far or close it is from other objects and may even later be able to adapt its behaviour accordingly. An example application that leveraged the location-aware paradigm was introduced by Hupfeld and Berge in their RAUM system ~\cite{hupfeld2000spatially}. The authors claim that information about the location of objects plays a more important role than their identities. They explain the essence of location by giving an example of people who prefer to communicate while standing in front of the person who participates in the conversation, rather than turning their backs on him. With the help of the concept of Ubiquitous Computing, the system presented in this paper uses the location information in order to select a communication partner.
\par Another issue related to the concept of Ubiquitous Computing is the scale, that is, systems of various sizes serve different purposes. In the context of our prototype, mascot, tablet, bench, and speakers are all in different sizes and, therefore, perform different tasks. Moreover, the size of objects is also reflected in its location, for example, a bench, compared to other devices has a larger size, which limits its location to one point, whereas the mascot which is a pocket-size phone allows changing the location depending on the location of its owner.

\subsection{Internet of Things}
\label{sec:Internet of Things}
The rapid development of electronics led to the emergence of the concept of “Internet of Things”. IoT can be both ubiquitous and non-ubiquitous technologies. Moreover, in the context of Ubiquitous Computing, IoT adds a new dimension to the interaction between objects: from any time, any place connectivity for everyone, we will have connectivity for anything ~\cite{tan2010future}. Thus, in comparison to Ubiquitous technologies, IoT focuses not only on the interaction between humans and devices but also between the devices themselves.
\par The idea of IoT was first proposed by Kevin Ashton in 1999 ~\cite{ashton2009internet}  by linking the idea of RFID (Radio Frequency Identification) to the topic of the Internet. We can characterize IoT as one big network where all devices can share information about their status with each other allowing to achieve deeper automation and integration within a system. In “Internet of Things: A Literature Review” paper ~\cite{madakam2015internet}, authors describe the genesis of the term “IoT” which help us to understand the general concept behind it and corresponding key technologies that it uses. They explained the concept by dividing the definition of IoT into two components: "Internet" (as a global system of interconnected computer networks that use the Internet protocol to serve users worldwide) and "Thing" (as real objects in the physical or material world). This explanation helps us to understand that inanimate objects such as light, speakers, chair, etc can communicate with other objects with the help of the Internet without any explicit human instructions. Thus, in our work, we use mascots, tablet, bench, and speakers as a representation of inanimate objects called “things”, which can interact with each other and send information over the local network. In addition, this paper provides key technology of IoT such as Radio Frequency Identification, Electronic Product Code, ZigBee, etc. From the technical point of view, RFID is primarily relevant to the unique identification of a “thing” in order to communicate with other objects. Moreover, ZigBee is widely used, short-range, low-rate wireless network technology, so in our prototype, the communication between mascot and bench is built with the help of Zigbee Lighting protocol. Additionally, an inexpensive radio technology Bluetooth Low Energy is also very useful for our research, which is used for proximity sensing. In addition to these technologies which are considered as a pillar for communication between objects, the more detailed description of their usage can be found in the Implementation chapter.
\par An example of research work in the area of IoT may be “Explorations on Reciprocal Interplay in Things Ecology”~\cite{chung2018explorations} where the authors are trying to stimulate scientists to a more detailed discussion on designing qualities of IoT devices. For that, Chung et al. conducted the HiddenLocal workshop (HWL) in order to explore and design IoT systems, where they take into account reciprocal interplay believing that it makes the design of IoT systems more dynamic. As a starting point, authors show 7 perceptual qualities as follows: focus the senses; show explorative behaviour; subtleness of movement; react to the external event; recognize explorative behaviour subject; reflex contextual noise; remember and anticipate perception over time. Authors believe that these perceptual qualities are a good approach for designed explorative features of devices and therefore for the things-to-things interaction.  

\subsection{Social Internet of Things}
\label{sec:Social Internet of Things}
According to the “The Internet of Things: A survey” paper ~\cite{atzori2010internet}, unfortunately, there are many research issues related to the IoT that require further research and need to be addressed. One of them is that people still cannot be sure about the privacy of the transferred data through IoT technologies. Another issue is network navigability which must ensure that the discovery of objects can be performed effectively and better reaction to the state changes of objects. Atzori et al. ~\cite{atzori2011siot} formalized a new paradigm of Social Internet of Things (SIoT) where the interaction among smart objects is based on the notion of Social Relationship of things rather than their owners. Thus, applying this concept to the IoT can lead to improving network navigability and the scalability. The architectural model of SIoT describes the establishment of the social relationships among objects in a fashion that is somewhat similar to the human social network relationship. 
\par Applying new paradigm to the concept of the IoT can lead to the following advantages: 
\begin{itemize}
  \item establishing the level of trustworthiness by leveraging relationship types and by supporting services usable among things that are “friends”; 
  \item improvements in network navigability; 
  \item a guarantee of the higher scalability and efficiency; ~\cite{atzori2012social};
\end{itemize}
\par By integrating social networking concept to the Internet of Things, intelligent things establish a connection with other peers in an autonomous by exploiting things' social relationships. An exemplary connection between smart things in our study may be Mascot-Mascot, Mascot-Tablet, and Mascot-Bench interactions. The application of the SIoT concept will help to accomplish complex tasks such as changing object behaviour according to the given information. Thus, with the help of the advantages provided by social networking principles the IoT evolves into the SIoT, in which social relationships can be established among the devices in order to advertise information about their current state and provide services to their peers.

\subsection{Autonomous interaction of things}
\label{sec:Autonomous interaction of things}
Most devices using the concept of IoT are designed to involve the user in the process where their actions trigger certain functions of a system in order to effects the behaviour of objects. This design contradicts to the concept of a fully automated system where objects can cooperate with each other beyond the control of a human. The following paper ~\cite{okada2016autonomous} , which is an inspiration for our work, introduces the design methodology to achieve a more autonomous system. The authors applied the concept of SIoT and consider objects as living beings which are able to communicate with others and exchange information autonomous. This approach allows objects to have their own social circle similar to human social network. This broadcast information calls certain functions that affect the behaviour of objects, thus, allowing objects to be aware of the status of other objects and the surrounding environment.
\par The concept of Social Things which is also an essential for our system helps the objects to know: own goals; what to do with the received information; and what actions need to be taken to achieve these goals. In our work, goals and the combination of actions that will be triggered depending on the received information, are all predefined.
\par As a case study, the authors developed a system with two devices. One of them is Mascot (which is a small keychain of three colours: red, green and blue) presented in the form of a personal object that the user can carry with him everywhere. Another device is a Bench with built-in lamps presented as a more static device for public use. Moreover, in the prototype described in that work, 2 scenarios are considered:
\begin{itemize}
  \item \textbf{Mascot - Mascot} interaction, as one mascot approaches another, they both start to blink where the intensity of blink depends on the distance objects are from each other.
  \item \textbf{Mascot - Bench} interaction, as Mascot goes close to the bench, the lights start to change their colors based on the color of mascot that is approaching.
\end{itemize}
The biggest contribution of this paper was to introduce the autonomously cooperative system where mascot and bench represented as a private and public thing. Moreover, authors also considered proximity-based cooperation: devices blink more often when approaching closer than 30 cm and and blink with less intensity when approaching more than 150 cm. By using the concept presented in this study (namely, autonomous interaction between objects/things achieved with the help of the SIoT concept), we are planning to expand the system by adding more objects. Afterwards, we are going to apply the theory of Proxemics and Personality Traits which will be covered in the following sections. In addition to the 2-categories that authors described in their paper (i.e. private and public "things") presented by 2 objects (mascot and bench, respectively), in section ~\ref{sec:The Theory of Proxemics} we are planning to look at more detailed divisions.

\subsection{The Theory of Proxemics}
\label{sec:The Theory of Proxemics}
Edward Hall ~\cite{hall1963system} ~\cite{hall1910hidden} conceptualized the idea of a personal space bubble by creating a whole system of notation in order to understand and record how people navigate shared space. He correlated physical distance to social distance.

According to these papers, Hall identified four distances which are measured horizontally:
\begin{itemize}
  \item \textbf{Intimate distance} which varies from 0 to 45 cm is a distance used for romantic partners and family members.
  \item \textbf{Personal distance} varies from 46 to 122 cm is a space bubble which allows your extended family members and close friends to enter this zone.
  \item \textbf{Social distance} varies from 122 to 370 cm is often used for acquaintances and colloquies
  \item \textbf{Public distance}, having a range of 370 cm and more, is often used in public speaking situation and with strangers you want to maintain your distance from.
\end{itemize}

He also analyzed vertical distances, for example, the difference in vertical distance between people can reflect the degree of dominance.

\par Nowadays, there are many studies in which Proxemics has been used to design interactions. For example, Jo Vermeulen et al. in their work ~\cite{vermeulen2015proxemic} used zones to interact with vertical interactive displays where they suggested floor display as an auxiliary device. The contributions of using the secondary display are the following: it provides peripheral information about tracking status of a user; it shows interaction zones; it invites the user to interact with the main display; it suggests possible interaction steps. This kind of floor visualization with continuous feedback about proximity gives the user more control over their interaction with the system.

\par Another example system that based its design on the theory of Proxemics is Remote Controls system introduced by Ledo et.al in their ~\cite{ledo2015proxemic} paper. Remote control devices were created in such a way that people could control appliances from a certain distance. However, with the increase in the number of home appliances, the number of remote controls also increasing. For this purpose, the universal remotes have been proposed providing a one-remote-to-many-appliances solution. Unfortunately, this design has setup issues and poorly adaptable interface. Authors of this paper presented proxemic-aware controls that utilize the spatial relationship between mobile devices owned by user and appliances surrounding it. With this system user can discover and select the devices within large ecologies of appliances, view their current status and control their features. Moreover, as a user moves closer or farther to a particular device, the interface adjusts accordingly. 
For example, in the initial state, the tablet screen visualizes icons representing the location of appliances at the edge of the screen, these icons are dynamically updated as he moves. Through spatial interactions, people can leverage mobile devices to discover and select appliances. This allows for situated interaction that balances simple and flexible control while seamlessly transitioning between different control interfaces. Ubicom may be a starting point for developing a new type of remote control interface within our increasingly complex world.

\par In addition, Ballendat, Nicolai Marquardt, and Saul Greenberg in their paper~\cite{ballendat2010proxemic} introduce proxemic-aware interactive media player system, where they consider information regarding nearby people and devices in order to mediate the interaction. They cover a small space Ubicomp environment considering the relationships of people to devices; devices to devices; and non-digital objects to people and devices. The system reacts to a person’s presence, distance and orientation regarding the display. Proxemic interaction also considers a person’s relationship to nearby objects. The authors propose different cases, for example, the video displaying on the screen pauses when a person is having a phone conversation or when he picks a magazine to read it. Another case is when a person enters the room, the screen shows a video title as additional information for him. Moreover, the video is paused when both people face away from the screen in order to start a conversation with each other. Furthermore, the system turns off when everyone has left the room. However, the authors also emphasize that one of the biggest unsolved problems in this area might be how the system can respond to the received information about proxemics because sometimes the devices can make a mistake by taking a certain action. In spite of all these problems, the authors, as well as we believe that proxemics will become an important factor in the embodiment of the interaction between social objects where they can meet the social expectations of people.

\par In our research, we are planning to extend the autonomous system, which was described in section ~\ref{sec:Autonomous interaction of things}, by categorizing devices according to the theory of Proxemics using only horizontal measurements. Our goal is to cover all four categories of the Theory of Proxemics which are represented by four artefacts (Mascot, Tablet, Bench and Speakers), thereby, by adding more case studies. Based on this theory, we are classifying our devices into following way: Mascot will belong to the intimate zone, Tablet to a Personal zone, Bench will be considered as a Social zone and Speakers as Public zone. In this way, these objects represent cooperation among these zones. According to Hall’s theory of Proxemics, the distance between people represents their relationship which affects the way how they interact with each other. Having understood how people use distance when interacting with each other, and then applying this concept to devices, based on this, we can come up with case studies and we hope that the interaction between devices will be more human-like. For example, the intimate zone has a very narrow range (i.e from 0 to 45 cm), and in the context of people, this distance used for romantic partners or for family members. Thus, by applying this zone in the context of a device-device relationship, we can come up with a device whose functionality is only visible for their owners (such as phone, watch or etc). Another example is Social zone, which is used for public speaking situations. In the context of SIoT, we can use Speakers as a representative of this zone. We suppose that the functionality of Speakers (i.e music play) will be available for everyone in the room. Thus, each of these devices is located at a certain distance from each other which represents the relationship between them. These relationships will help us to conceptualize their interactions, come up with case study and possible actions.

\subsection{The Social Structure of things}
\label{sec:The Social Structure of things}
Given that the Theory of Proxemics correlates the physical distance with the relationship, we need to explore what kind of relationship these objects can establish. Atzori et al in their “SIoT: Giving a Social Structure to the IoT” paper ~\cite{atzori2011siot} discussed several types of social relationship according to which things have to interact with each other. They described common forms of sociality that objects may use in their relations. According to the paper, objects can establish the following types of relationships: \textit{parental object relationship} built by the same manufacturer which will not change over time and is only updated by disruption of a given object; \textit{co-location and co-work object relationship} with a frequency of changes based on the physical location and amount of the interaction; \textit{ownership object relationship} when the same user owned certain objects; \textit{social object relationship} when authorized objects are able to share certain information about themselves with each other in an autonomous way. Meaning that the devices with similar profiles are able to share their best experiences in order to help other devices to solve problems with whom the first have already encountered. This type of relationship provides a distributed solution that is expected to be efficient. Thus, the relationship between social objects in the context of SIoT needs more focus and a further examination. Author of the following work ~\cite{fiske1992four} provided helpful hints in this direction, which was subsequently applied in the context of the SIoT in the work of ~\cite{pintus2015humanizing}. We will describe this work in more details in the following subsection.

\subsection{Social Patterns in the Humanized IoT Domain}
\label{sec:Social Patterns in the Humanized IoT Domain}
Considering the type of relationship that social objects may have, there is a need to find out which actions can be established within this relationship. Alan Fiske in his ~\cite{fiske1992four} paper defined Rational Model Theory where according to that model people interact and establish goals and actions with each other. The following paper ~\cite{pintus2015humanizing} defined a new pattern based on the Alan Fiske “relational model” theory ~\cite{fiske1992four}. Friske’s four elementary forms of sociality projected the idea of Humanized Social Internet of Things which includes both things and people in its foundation. According to this model, people establish social actions with each other through "things". They described 4 forms which are Communal Sharing, Authority Ranking, Equality Matching, and Market Pricing where:
\begin{itemize}
  \item  \textbf{Communal Sharing} represents a trusted sharing of things in which people from the same community are allowed to use each other's devices which means that other people have the same level of control as the owner of this device.
  \item  \textbf{Authority Ranking} represents the hierarchy which restricts people who are not owners of these devices to do certain actions. 
  \item  \textbf{Equality Matching} allows people to share smart things in order to reach common goals.
  \item  \textbf{Market Pricing} provides a transaction-based model, where people exchange devices, and money is often part of the deal. From the IoT perspective, any kind of physical objects or proprietary software might serve as a good example.
\end{itemize}

\par In our study, we would like to combine two models and apply in the context of SIoT, where social things establish social actions with each other through functionalities / services that they offer. Since, tablet, bench and speakers are static devices and mascots are both dynamic and ubiquitous, we will consider the relationship between mascots in relation to the functionality of static objects.
\par In one case study, when several mascots interact with objects (i.e with each other, tablet and/or bench), the server (which, in our prototype, is the main decision point) chooses the most active mascot and then speakers play music for it. For example, the relationship between these mascots in relation to the use of the Speakers functionality will be Communal Sharing Relationship. Since each mascot will have equal access to the functionality of a Speaker, there will be no queues in our model during the use of this static device. Each Mascot will have the same level of control/permission over Speakers. The reason why in this case we do not use Equality Matching instead of Communal Sharing (CS) relationship, is because the concept of Equality Matching (EM) is order preserved. That is, if the music already played for Mascot-1, it will have to wait for his next turn until the music will be played for all other mascots. This approach or architecture is not considered by our system since we want to achieve more intuitive and more natural interaction between devices as it would be between people. In case of CS, if the music is already played for the Mascot-1, as long as this mascot is chosen as the most active one, it can be played for it again, out of turn.
\par In another case study, when Mascot approaches another Mascot, they start to vibrate. So there is a relationship between Mascots through their functionality (i.e vibration). At this case, we plan to apply Authority Ranking Relationship. Meaning that the person who owned the mascot gives permission to other Mascot owners to change its state, but it does not share authority over it, thereby, ensuring that authority cannot be changed by people who do not own the particular mascot. In other words, when someone reaches you in less than 45 cm (which based on Proxemics theory is considered as an intimate zone), your mascot will change its status by vibrating. However, the change in the state of your device can be limited and controlled by moving from one place to another, thereby not allowing the other person, who also holds his own mascot, to enter your intimate distance. Thus, there is a hierarchy between mascots, and as an owner of a mascot if you do not want your mascot to be vibrated, you can just stay far from other Mascots by that you limit permission to have access to the functionality of your Mascot.

\par Such a distribution of devices (namely, classification of fixed devices to “Communal Sharing”, and mascots, which are more dynamic, to “Authority Ranking” category) helps us better define goals and understand the behaviour of each object during the interaction with each other. Moreover, applying this kind of Relational Model may help us to conceptualise the behaviour and limits the possible actions of devices in the same situation.

\subsection{Interaction design for SIoT}
\label{sec:Interaction design for SIoT}

\par The following paper ~\cite{soro2018social} motivated us to apply a concept of personality in the context of social devices. An example use of Personality as a method to design an interactive object’s behavior was proposed in the “Designing the Behaviour of Interactive Objects” paper ~\cite{spadafora2016designing}. The authors came to the conclusion that in order to design a more stable and understandable for user the behavior of a device, it is necessary to add inner logic to which we can refer. Marco et al. proposed to apply the concept of metaphor, which represents human stereotypes of personality in order to visualise the inner logic. Their system was based on a Big Five Personality Traits model, and thus, by assigning these personalities to objects, users could more easily describe its behavior. The authors believe that stereotypes and metaphors are simplified descriptions of being and behavior, and thus making it an ideal method for displaying the sustainable behavior of a smart object. During the research, they used robotic sofa as a case study and tried to analyze how users perceive the consistency of its behavior. The use of the Personality model in the device design process, helps a user to create a mental model of how an autonomous sofa-bot will act in the future.

Having a system, where tablet, bench and speakers are considered as static objects, whereas mascots are dynamic, we can apply the Personality Model. And since only mascots are a major factor affecting the environment (for example, if mascot come close to the bench, it changes the light color), so we decided to assign a personality to dynamic objects (meaning to each mascot).

In our prototype we assume that this approach may help user to better understand the mascot behavior. Knowing which goals, intentions this object follows may help users to understand the behavior and the reason of certain decisions of a mascot. We believe that these goals and intentions set certain boundaries in behavior of social devices. And since the behavior as a whole consists of actions, we suppose that personality model sets implicit instructions to the devices. Meaning that, instead of the user giving explicit instructions to the object, the object with the assigned personality makes decisions autonomously. Thus, the device with assigned personality can help users who knows the definition of that personality to understand the system behaviour at least in an intuitive way.
This concept may give the system a more understandable and consistent behavior, and to user a better awareness of object functionality
In the following subsection we will describe the Personality Model in more details.

\subsubsection{Definition of Personality Traits}
\label{sec:Definition of Personality Traits}

Personality is important in human relationship, so we hope that it also may be important for device relationship. In order to assign a personality to each mascot, we first need to give a definition of personality in the context of social devices. We can try on an intuitively explain the meaning of a person’s personality trait, unfortunately, it is hard to apply it in the context of SIoT. For that, we need a generally accepted model, and we decided to use Big Five Personality Traits (aka OCEAN) for the description of each personality. We hope that providing a description of personality will help to define goals and more targeted actions which in turn will lead the system to more stable behavior. The following book ~\cite{matthews2003personality} gives a good introduction to the personality types describing possible existing personality models. One of the models that can be used was introduced by Costa and McCrae’s five-factor model which is also known as Big Five Personality Traits and the OCEAN model. Moreover, their concept formed a basis for the widely used NEO-Personality Inventory-Revised (NEO-PI-R) measurement scale.

\par The OCEAN model consists of the following features: Openness to experience, Conscientiousness, Extroversion, Agreeableness and Neuroticism. The authors lists the facets associated with each of these five domains:

\begin{itemize}
  \item  \textbf{Openness to experience}: creativity, innovative quality; quick receptivity to new and abstract ideas, high intelligence and openness to novelty;
  \item  \textbf{Conscientiousness}: organized, well-prepared, discipline, likes planned action more than spontaneity, more focused.
  \item  \textbf{Extroversion}: energetic, assertive personality, like to be the center of attention, like to dominate, feel comfortable around people;
  \item  \textbf{Agreeableness}: friendliness, compassion for other people, interested in people, sympathize with the feelings of others, soft-hearted;
  \item  \textbf{Neuroticism}:  irritability, more hostile towards others, most often feel anxiety when they are surrounded by others, frequent mood swings, emotionally unstable;
  \end{itemize}
  
\par We assigned a set of personalities to dynamic objects each of which is described in the above-mentioned list. In addition, in spite of the fact that people usually have a combination of these five traits, we are going to consider only extreme cases. If the mascot is assigned extravert personality, we are planning to consider only this trait, by making this device, for example, highly assertive, and extremely energetic.  


\clearpage
\begingroup
  \pagestyle{plain}
  \bibliographystyle{unsrt}
   \phantomsection
	\addcontentsline{toc}{section}{Bibliography}
	\bibliography{references}\emph{}	
	\vfill
  \clearpage
\endgroup 
\end{document}